% ============================================================================
%  sections/conclusion.tex
% ============================================================================
\chapter*{Conclusion générale}
\addcontentsline{toc}{chapter}{Conclusion générale}

Le stage effectué au sein de l'ENSA Béni Mellal nous a permis de concevoir et développer une plateforme web complète de présélection des candidats aux filières d'ingénieur. Ce projet a couvert l'ensemble du cycle de développement logiciel : analyse des besoins, conception architecturale, implémentation Full Stack et tests.

La plateforme développée répond aux objectifs fixés : automatisation du traitement des dossiers via un moteur de règles configurable à 7 types, scoring académique pondéré par filière avec fuzzy matching des matières, gestion des épreuves écrites avec import Excel, notifications en temps réel via WebSocket, et traçabilité complète des décisions.

Sur le plan technique, l'architecture Django REST Framework + React 18 s'est révélée robuste et performante. Les 5 applications Django assurent une séparation claire des responsabilités, tandis que le frontend React offre une expérience utilisateur fluide et responsive. La stratégie de tests a permis d'atteindre un taux de réussite de 100\% sur 61 tests avec une couverture de code de 87\%.

Ce stage a été une expérience enrichissante qui m'a permis de consolider mes compétences en développement Full Stack, en architecture logicielle et en gestion de projet agile.

\section*{Perspectives d'amélioration}
\addcontentsline{toc}{section}{Perspectives d'amélioration}

\begin{itemize}[leftmargin=2cm]
  \item Intégration avec la plateforme nationale Massar via convention ministérielle pour la vérification automatique des CNE.
  \item Développement d'une application mobile (Flutter/React Native) pour le suivi des candidatures en mobilité.
  \item Tableau de bord analytique avancé avec indicateurs de performance détaillés.
  \item Intelligence artificielle pour la détection de fraude documentaire avancée.
  \item Déploiement en production avec conteneurisation Docker et CI/CD.
\end{itemize}
